% source: eb1c5102987a86203e9a24cce85643a80d53317f
\chapter{ML-DSA（FIPS 204）的完整实现}
\chaptermark{ML-DSA（FIPS 204）}
\label{ch:mldsa}

\section{学习目标}

本书前一部分一路推进到 FIPS 203，即用于\emph{加密与密钥交换}的正式 ML-KEM 标准。现在我们转向公钥密码的另一半——\emph{数字签名}，并进入三个签名标准中的第一个：FIPS 204，即 \textbf{ML-DSA}，它更为人熟知的是其提交时的名称 Dilithium~\cite{fips204,dilithium2018}。

学完本章后，你应当能够：

\begin{enumerate}
  \item 说明签名方案做什么，以及它与 KEM 有何不同；
  \item 理解 ML-DSA 的密钥生成流程，并看到它几乎原封不动地复用了 Kyber 的模格机制；
  \item 说明两个困难问题——模 LWE 与模 SIS——以及它们各自保护方案的哪一部分；
  \item 逐阶段理解签名过程，并明白其中为何包含一个拒绝（“中止”）循环；
  \item 把拒绝的思想与 LWE 各章所研究的噪声作用联系起来。
\end{enumerate}

\section{回顾：签名与 ML-KEM 有何不同}

ML-KEM 回答的问题是：\emph{双方如何就一个共享秘密达成一致？}签名回答的是另一个问题：\emph{验证者如何知道一条消息确实来自秘密钥的持有者，且未被改动？}三个算法在形式上与 KEM 相对应，但服务于这一新目标：

\begin{itemize}
  \item \textbf{KeyGen} $\to$ 公开的验证密钥 $pk$ 与秘密的签名密钥 $sk$；
  \item \textbf{Sign}$(sk, M)$ $\to$ 消息 $M$ 上的签名 $\sigma$；
  \item \textbf{Verify}$(pk, M, \sigma)$ $\to$ \emph{接受}或\emph{拒绝}。
\end{itemize}

安全目标是\emph{不可伪造性}：即使看过许多有效签名，没有 $sk$ 的人也无法对一条新消息生成有效签名。值得注意的是，ML-DSA 正是在我们为 Kyber 构建的\emph{完全相同}的环 $R_q = \mathbb{Z}_q[x]/(x^{256}+1)$ 和模结构上实现这一目标的。只有模数改变了（$q = 8380417$），而签名过程则是真正全新的。

\section{两个困难问题}

ML-DSA 依赖于第~\ref{ch:mlwe}~章模设定下的两个格问题。

\begin{itemize}
  \item \textbf{模 LWE}，我们已经熟悉：给定 $(A,\ \mathbf{t}=A\mathbf{s}_1+\mathbf{s}_2)$，其中 $\mathbf{s}_1,\mathbf{s}_2$ 较短，恢复秘密是困难的。它把签名密钥隐藏在公钥之中。
  \item \textbf{模 SIS}（短整数解，Short Integer Solution）：第~\ref{sec:sis}~节引入的 SIS 问题的模格形式。给定随机的 $A$，找到满足 $A\mathbf{z}=\mathbf{0}\pmod q$ 的短非零向量 $\mathbf{z}$ 是困难的。它阻止伪造：伪造者必须给出一个满足某个公开方程的短向量。
\end{itemize}

SIS 是 LWE 的镜像。LWE 把秘密隐藏在噪声之中；SIS 则要求你为一个随机方程组\emph{找到}一个短解。二者困难的几何原因相同——在高维格中很难找到短向量。

\section{FIPS 204 中的真实输入}

与 FIPS 203 完全一样，密钥生成并不从采样矩阵开始，而是从单个短种子开始：
\[
\xi \in \{0,1\}^{256}.
\]
其余一切——公开矩阵、两个秘密向量以及内部的 PRF 密钥——都由 $\xi$ 确定性地扩展得到。因此，与 ML-KEM 一样，密钥生成是一个确定性函数 $\mathrm{KeyGen}(\xi)$。

\section{阶段 1：扩展种子}

种子经 $H$（基于 SHAKE-256 构建）哈希后分为三部分：
\[
(\rho,\ \rho',\ K) = H(\xi).
\]
它们的作用是：
\begin{itemize}
  \item $\rho$——公开，扩展为矩阵 $A$（它将存储在公钥中）；
  \item $\rho'$——秘密，作为采样 $\mathbf{s}_1$ 与 $\mathbf{s}_2$ 的种子；
  \item $K$——秘密，一个 PRF 密钥，稍后用于派生每次签名的随机性。
\end{itemize}

\section{阶段 2：生成矩阵}

与 Kyber 一样，矩阵从不存储，而是按需从其种子重新生成：
\[
A = \textsc{ExpandA}(\rho) \in R_q^{k\times \ell}.
\]
维数 $(k,\ell)$ 正是区分三个安全级别的参数。

\section{阶段 3：采样秘密向量}

两个短秘密由 $\rho'$ 扩展得到：
\[
(\mathbf{s}_1,\ \mathbf{s}_2) = \textsc{ExpandS}(\rho'),\qquad \mathbf{s}_1\in R_q^\ell,\ \mathbf{s}_2\in R_q^k,
\]
其系数位于较小的范围 $[-\eta,\eta]$ 内。它们扮演 LWE 秘密与误差的角色。

\section{阶段 4：计算并拆分公开值}

公开值就是熟悉的模 LWE 方程：
\[
\mathbf{t} = A\mathbf{s}_1 + \mathbf{s}_2.
\]
为缩小公钥，$\mathbf{t}$ 被拆分为高位部分 $\mathbf{t}_1$（公开）和低位部分 $\mathbf{t}_0$（保存在秘密钥中）：
\[
\mathbf{t} = \mathbf{t}_1\cdot 2^{d} + \mathbf{t}_0.
\]
公钥为 $(\rho, \mathbf{t}_1)$；秘密钥则打包了 $(\rho, K, \mathbf{t}_0, \mathbf{s}_1, \mathbf{s}_2)$，以及签名时使用的哈希值 $tr = H(pk)$。

\begin{figure}[H]
\centering
\begin{tikzpicture}[every node/.style={font=\scriptsize}]
  \node[flownode] (xi) at (5,7.4) {seed $\xi \in \{0,1\}^{256}$};
  \node[flownode] (h)  at (5,6.2) {$H(\xi)$};
  \node[flownode] (rho)   at (1.3,5.0) {$\rho$};
  \node[flownode] (rhop)  at (5,5.0) {$\rho'$};
  \node[flownode] (K)     at (8.7,5.0) {$K$ (PRF key)};
  \node[flownode, font=\scriptsize] (A)     at (1.5,3.6) {$A=\textsc{ExpandA}(\rho)$};
  \node[flownode, font=\scriptsize] (s)     at (5.9,3.6) {$(\mathbf{s}_1,\mathbf{s}_2)=\textsc{ExpandS}(\rho')$};
  \node[flownode] (t)     at (3.6,2.4) {$\mathbf{t}=A\mathbf{s}_1+\mathbf{s}_2$};
  \node[flownode] (split) at (3.6,1.3) {split $\mathbf{t}=\mathbf{t}_1 2^d+\mathbf{t}_0$};
  \node[keynode]  (pk)    at (1.7,-0.1) {public key $(\rho,\mathbf{t}_1)$};
  \node[keynode]  (sk)    at (7.0,-0.1) {secret key\\[-2pt]$(\rho,K,tr,\mathbf{t}_0,\mathbf{s}_1,\mathbf{s}_2)$};
  \draw[flowarrow] (xi) -- (h);
  \draw[flowarrow] (h) -- (rho);
  \draw[flowarrow] (h) -- (rhop);
  \draw[flowarrow] (h) -- (K);
  \draw[flowarrow] (rho) -- (A);
  \draw[flowarrow] (rhop) -- (s);
  \draw[flowarrow] (A) -- (t);
  \draw[flowarrow] (s) -- (t);
  \draw[flowarrow] (t) -- (split);
  \draw[flowarrow] (split) -- (pk);
  \draw[flowarrow] (split) -- (sk);
  \draw[flowarrow] (K) to[out=270,in=0] (sk.east);
\end{tikzpicture}
\caption{ML-DSA 密钥生成，采用与 FIPS 203 相同的确定性扩展风格。单个种子 $\xi$ 分散生成矩阵、两个秘密以及 PRF 密钥；公钥只存储 $(\rho,\mathbf{t}_1)$，由它总能重新生成 $A$。}
\label{fig:mldsa-keygen}
\end{figure}

\begin{algorithm}[H]
\caption{ML-DSA.KeyGen（教学形式）}
\label{alg:mldsa-keygen}
\begin{algorithmic}[1]
\Require 种子 $\xi \in \{0,1\}^{256}$；参数 $(k,\ell,\eta,d)$
\Ensure 公钥 $(\rho,\mathbf{t}_1)$；秘密钥 $(\rho,K,tr,\mathbf{t}_0,\mathbf{s}_1,\mathbf{s}_2)$
\State $(\rho,\rho',K) \gets H(\xi)$ \Comment{阶段 1}
\State $A \gets \textsc{ExpandA}(\rho)$ \Comment{阶段 2}
\State $(\mathbf{s}_1,\mathbf{s}_2) \gets \textsc{ExpandS}(\rho')$ \Comment{阶段 3}
\State $\mathbf{t} \gets A\mathbf{s}_1 + \mathbf{s}_2$ \Comment{阶段 4}
\State $(\mathbf{t}_1,\mathbf{t}_0) \gets \textsc{Power2Round}(\mathbf{t}, d)$
\State $tr \gets H(\rho\,\|\,\mathbf{t}_1)$
\State \Return $\big((\rho,\mathbf{t}_1),\ (\rho,K,tr,\mathbf{t}_0,\mathbf{s}_1,\mathbf{s}_2)\big)$
\end{algorithmic}
\end{algorithm}

\section{插曲：Fiat--Shamir 变换}
\label{sec:fiat-shamir}

在深入 ML-DSA 的签名之前，值得先把它所用的通用技术单独拿出来讨论，因为这是一个基本构件，其应用远不止这一个方案。

\textbf{身份识别方案}（identification scheme）让证明者在不泄露 $s$ 的前提下，使验证者相信它知道秘密 $s$（其公开像为 $t$）。这是一段三轮对话：证明者发送\emph{承诺} $w$；验证者回复随机\emph{挑战} $c$；证明者发送只有知道 $s$ 的人才能给出的\emph{响应} $z$，验证者检查 $(w,c,z,t)$ 之间的某个关系。如果响应经过恰当的随机化，则整段记录不泄露关于 $s$ 的任何信息。

\textbf{Fiat--Shamir 变换}~\cite{fiatshamir1986}通过\emph{去掉验证者}，把这一交互式协议变为签名。证明者不再等待随机挑战，而是自己把挑战计算为承诺与消息一起的哈希值：
\[
c = H(w \,\|\, m).
\]
由于 $H$ 的行为如同一个不可预测的随机函数，证明者无法把 $c$ 引向对自己有利的方向；又由于 $c$ 依赖于 $m$，整段记录便与该消息绑定——它\emph{就是} $m$ 上的一个签名。

\begin{figure}[H]
\centering
\begin{tikzpicture}[font=\scriptsize]
  \node[font=\small\bfseries] at (1.75,3.5) {Interactive identification};
  \node[seqparty] (P) at (0,2.9) {Prover};
  \node[seqparty] (V) at (3.5,2.9) {Verifier};
  \draw[seqlife] (0,2.6) -- (0,0.1);
  \draw[seqlife] (3.5,2.6) -- (3.5,0.1);
  \draw[seqmsg] (0,2.0) -- node[above]{commitment $w$} (3.5,2.0);
  \draw[seqmsg] (3.5,1.35) -- node[above]{random $c$} (0,1.35);
  \draw[seqmsg] (0,0.7) -- node[above]{response $z$} (3.5,0.7);
  \node[seqnote, left] at (3.35,0.3) {checks $(w,c,z,t)$};
  \node[pqcorange, font=\Large] at (5.1,1.6) {$\Rightarrow$};
  \node[font=\small\bfseries] at (8.2,3.5) {Non-interactive (Fiat--Shamir)};
  \node[flownode, align=center] (P2) at (8.2,2.3) {Prover computes $w$,\\then $c=H(w\,\|\,m)$ itself};
  \node[keynode] (sig) at (8.2,1.0) {signature $\sigma=(c,z)$};
  \draw[flowarrow] (P2) -- (sig);
\end{tikzpicture}
\caption{Fiat--Shamir 变换。左侧是一个交互式身份识别协议：共三条消息，由验证者提供新鲜的随机挑战 $c$。右侧是同一段记录的非交互化版本：由证明者自己计算 $c=H(w\,\|\,m)$——哈希扮演验证者的角色，而将其与消息 $m$ 绑定，便使这段记录成为一个签名。许多签名方案都以这一配方为基础；ML-DSA 用模格将其实例化。}
\label{fig:fiat-shamir}
\end{figure}

\section{签名的思想：承诺、挑战、响应}

签名恰好实例化了上述身份识别记录，由模格对象分别扮演承诺、挑战和响应的角色：

\begin{enumerate}
  \item \textbf{承诺。}采样一个随机的短掩码向量 $\mathbf{y}$，并构造 $\mathbf{w} = A\mathbf{y}$。
  \item \textbf{挑战。}先构造消息代表值 $\mu=H(tr\,\|\,M')$，其中 $M'$ 包含任何上下文或预哈希格式。然后计算 $\tilde c=H(\mu\,\|\,\mathbf{w}_1)$，其中 $\mathbf{w}_1=\textsc{HighBits}(\mathbf{w})$，并导出 $c=\textsc{SampleInBall}(\tilde c)$。
  \item \textbf{响应。}以 $\mathbf{z} = \mathbf{y} + c\,\mathbf{s}_1$ 作答。
\end{enumerate}

于是验证者可以重新计算 $A\mathbf{z} - c\,\mathbf{t}$，并检查它在高位上与 $\mathbf{w}$ 一致，因为
\[
A\mathbf{z} - c\,\mathbf{t} = A(\mathbf{y}+c\,\mathbf{s}_1) - c(A\mathbf{s}_1+\mathbf{s}_2) = A\mathbf{y} - c\,\mathbf{s}_2 \approx \mathbf{w}.
\]
剩余的 $c\,\mathbf{s}_2$ 很小，\textsc{HighBits} 的舍入会将其吸收。用哈希 $H$ 取代交互式验证者，正是 Fiat--Shamir 变换：哈希充当一个不可预测、不可操纵的挑战者。

\section{为什么签名有时必须中止}

响应 $\mathbf{z} = \mathbf{y} + c\,\mathbf{s}_1$ 直接携带了秘密 $\mathbf{s}_1$。如果掩码 $\mathbf{y}$ 不够大，不足以淹没 $c\,\mathbf{s}_1$，那么 $\mathbf{z}$ 就会泄露信息——在许多签名上取平均后，掩码相互抵消，秘密便显露出来。这是换了装扮的 LWE 教训：\emph{掩码必须完全隐藏秘密}。

ML-DSA 的解决办法是\textbf{拒绝采样}（rejection sampling），也称为\emph{带中止的} Fiat--Shamir。构造出 $\mathbf{z}$ 后，签名者检查它是否落在一个安全区域内，在该区域内其分布与 $\mathbf{s}_1$ 无关。如果不是，就丢弃一切，重新抽取新的 $\mathbf{y}$ 并重试。只有不泄露信息的 $\mathbf{z}$ 才会被发布。

\begin{figure}[H]
\centering
\begin{tikzpicture}[node distance=5.5mm and 11mm]
  \node[flownode] (y) {sample mask $\mathbf{y}$};
  \node[flownode, right=of y] (w) {$\mathbf{w}=A\mathbf{y}$};
  \node[flownode, right=of w] (c) {$\tilde c=H(\mu\,\|\,\mathbf{w}_1)$};
  \node[flownode, below=of c] (z) {$\mathbf{z}=\mathbf{y}+c\,\mathbf{s}_1$};
  \node[diamond, draw=pqcorange, thick, fill=white, align=center,
        inner sep=1pt, font=\scriptsize, left=of z] (chk) {$\mathbf{z}$\\safe?};
  \node[keynode, left=of chk] (out) {output\\[-2pt]\scriptsize $\sigma=(c,\mathbf{z},\mathbf{h})$};
  \draw[flowarrow] (y) -- (w);
  \draw[flowarrow] (w) -- (c);
  \draw[flowarrow] (c) -- (z);
  \draw[flowarrow] (z) -- (chk);
  \draw[flowarrow] (chk) -- node[above,font=\scriptsize,pqcteal]{yes} (out);
  \draw[flowarrow] (chk.north) |- ([yshift=3.5mm]y.north)
        node[near end,above,font=\scriptsize,pqcorange]{no: restart} -- (y.north);
\end{tikzpicture}
\caption{ML-DSA 签名循环。每一轮都抽取新的掩码 $\mathbf{y}$；只有当安全检查确认响应 $\mathbf{z}$ 的分布与 $\mathbf{s}_1$ 无关时，才会发布它。平均只需少数几次迭代。提示 $\mathbf{h}$ 使验证者即使在低位部分 $\mathbf{t}_0$ 被舍弃的情况下，也能恢复高位。}
\label{fig:mldsa-abort}
\end{figure}

\section{签名与验证}

把各阶段合在一起，就得到签名算法。安全区域由两个界刻画：$\|\mathbf{z}\|_\infty$ 必须保持在 $\gamma_1-\beta$ 以下，且 $\mathbf{w}-c\mathbf{s}_2$ 的低位必须保持在 $\gamma_2-\beta$ 以下。提示 $\mathbf{h}$ 用于弥补被省略的 $\mathbf{t}_0$。

\begin{algorithm}[H]
\caption{ML-DSA.Sign（教学形式）}
\label{alg:mldsa-sign}
\begin{algorithmic}[1]
\Require 秘密钥 $(\rho,K,tr,\mathbf{t}_0,\mathbf{s}_1,\mathbf{s}_2)$，消息 $M$，32 字节值 $rnd$
\Ensure 签名 $\sigma=(c,\mathbf{z},\mathbf{h})$
\State $A \gets \textsc{ExpandA}(\rho)$;\quad $\mu \gets H(tr\,\|\,M')$
\State $\kappa \gets 0$;\quad $\rho'' \gets H(K\,\|\,rnd\,\|\,\mu)$ \Comment{$rnd=0$ 对应确定性选项}
\Loop
  \State $\mathbf{y} \gets \textsc{ExpandMask}(\rho'', \kappa)$;\quad $\kappa \gets \kappa+\ell$ \Comment{每轮使用新掩码}
  \State $\mathbf{w} \gets A\mathbf{y}$;\quad $\mathbf{w}_1 \gets \textsc{HighBits}(\mathbf{w})$
  \State $\tilde c \gets H(\mu\,\|\,\mathbf{w}_1)$;\quad $c\gets\textsc{SampleInBall}(\tilde c)$
  \State $\mathbf{z} \gets \mathbf{y} + c\,\mathbf{s}_1$
  \If{$\|\mathbf{z}\|_\infty \ge \gamma_1-\beta$ \textbf{ or } $\|\textsc{LowBits}(\mathbf{w}-c\mathbf{s}_2)\|_\infty \ge \gamma_2-\beta$}
    \State \textbf{continue} \Comment{不安全：丢弃并重新采样}
  \EndIf
  \State $\mathbf{h} \gets \textsc{MakeHint}(c\mathbf{t}_0,\ \mathbf{w}-c\mathbf{s}_2+c\mathbf{t}_0)$
  \State \Return $(c,\mathbf{z},\mathbf{h})$
\EndLoop
\end{algorithmic}
\end{algorithm}

验证过程由公开数据重新计算承诺，检查挑战哈希是否一致，同时确认 $\mathbf{z}$ 确实是短的——正是这一短向量检查让伪造者面对模 SIS 问题。

\begin{algorithm}[H]
\caption{ML-DSA.Verify（教学形式）}
\label{alg:mldsa-verify}
\begin{algorithmic}[1]
\Require 公钥 $(\rho,\mathbf{t}_1)$，消息 $M$，签名 $(c,\mathbf{z},\mathbf{h})$
\Ensure \emph{接受}或\emph{拒绝}
\State $A \gets \textsc{ExpandA}(\rho)$;\quad $\mu \gets H(H(\rho\,\|\,\mathbf{t}_1)\,\|\,M)$
\State $\mathbf{w}_1' \gets \textsc{UseHint}\big(\mathbf{h},\ A\mathbf{z} - c\,\mathbf{t}_1\cdot 2^d\big)$
\State $c' \gets H(\mu\,\|\,\mathbf{w}_1')$
\If{$c'=c$ \textbf{ and } $\|\mathbf{z}\|_\infty < \gamma_1-\beta$}
  \State \Return \emph{接受}
\Else
  \State \Return \emph{拒绝}
\EndIf
\end{algorithmic}
\end{algorithm}

由于 $A\mathbf{z} - c\,\mathbf{t} = A\mathbf{y} - c\mathbf{s}_2$ 在高位上与 $\mathbf{w}=A\mathbf{y}$ 一致，而提示又修补了只公开 $\mathbf{t}_1$ 所留下的差距，因此当且仅当签名真实时，双方才会哈希出相同的挑战。

\section{伪造者为何无路可走}

没有 $sk$ 的攻击者可以任意选择 $\mathbf{z}$ 和 $c$，但要通过验证，就必须满足 $c = H(\mu\,\|\,\textsc{HighBits}(A\mathbf{z}-c\mathbf{t}))$，\emph{并且}保持 $\mathbf{z}$ 是短的。由于 $H$ 的行为如同随机函数，其输出无法被操控；唯一的途径是找到一个与该方程相容的真正短的 $\mathbf{z}$——这就是一个模 SIS 实例。而攻击者也无法转而从公开数据中获知 $\mathbf{s}_1$，因为那是模 LWE 问题。两个困难问题同时堵住了两扇门。

\section{SageMath 实验}

核心的抵消关系很容易在缩小版中验证。使用小参数，采样一个密钥，构造一个加了掩码的响应，并在发布任何签名之前就确认验证恒等式成立。

\begin{lstlisting}[language=Python]
# Toy sizes; real ML-DSA uses n=256, q=8380417.
from random import Random
from sage.all import GF, Matrix, PolynomialRing, vector

q = 8380417
R.<x> = PolynomialRing(GF(q))
Rq = R.quotient(x^8 + 1)  # toy ring
rng = Random("mldsa-cancellation")

def small_poly():
    return Rq([rng.choice([-1, 0, 1]) for _ in range(8)])

def uniform_poly():
    return Rq([rng.randrange(q) for _ in range(8)])

A = Matrix(Rq, 2, 2, [uniform_poly() for _ in range(4)])
s1 = vector(Rq, [small_poly(), small_poly()])
s2 = vector(Rq, [small_poly(), small_poly()])
t = A * s1 + s2
y = vector(Rq, [uniform_poly(), uniform_poly()])
c = small_poly()
# For a mask y and challenge c, check the verifier's identity:
#   A*z - c*t  ==  A*y - c*s2   (should hold exactly)
z = y + c*s1
assert A*z - c*t == A*y - c*s2
\end{lstlisting}

然后用若干个掩码重复实验，统计 $\|\mathbf{z}\|_\infty$ 超出界的频率：这恰好就是拒绝率，它说明了为什么一次签名需要少量且次数不定的尝试。

\section{完整的 SageMath 实现}
\label{sec:fips204-impl}

上面的抵消检查只是一个片段。配套文件
\sagefile{fips204\_mldsa.sage} 则是整个标准：FIPS~204 中全部四十九个编号算法，每个都实现为独立的函数并注明其算法编号，覆盖 ML-DSA-44、65 和 87。其中包括本章当作黑盒处理的部分——\textsc{ExpandA}、\textsc{ExpandS}、
\textsc{ExpandMask}、\textsc{SampleInBall}、\textsc{Power2Round}、
\textsc{Decompose}、\textsc{MakeHint}、\textsc{UseHint}——以及位打包、提示编码、预哈希变体 \textsc{HashML-DSA} 和
Montgomery 约简。

\playground{ML-DSA-44、65 和 87 三个参数集上的密钥生成、签名与验证}

与 FIPS~203 一样，代数运算由 Sage 承载。不过这里标准细心地区分了环 $R_q$ 与系数范围受限的整系数多项式环 $R$，代码也同样如此：模 $q$ 约简的量（秘密 $\mathbf{s}_1,\mathbf{s}_2$、掩码 $\mathbf{y}$、响应 $\mathbf{z}$）位于 Sage 商环
$R_q = \mathbb{F}_{8380417}[X]/(X^{256}+1)$ 中，而 $\mathbf{t}_1$、
$\mathbf{w}_1$、$\mathbf{r}_0$ 以及提示 $\mathbf{h}$ 则保持为整数列表。
NTT 域 $T_q$ 是直积 $(\mathbb{Z}_q)^{256}$，因此
\textsc{MultiplyNTT}（算法 45）实际上就是逐坐标乘积。

于是，算法~\ref{alg:mldsa-sign} 的拒绝循环读起来几乎与标准中的写法完全一致：

\begin{lstlisting}[language=Python]
kappa = 0
while True:
    y  = self.ExpandMask(rhopp, kappa)
    yh = [NTT(Rq(p)) for p in y]
    w  = [NTTInverse(a) for a in MatrixVectorNTT(Ah, yh)]
    w1 = HighBitsVec(w, self.gamma2)

    ctilde = H(mu + self.w1Encode(w1), self.lam // 4)
    c      = self.SampleInBall(ctilde)
    ch     = NTT(Rq(c))

    cs1  = [NTTInverse(MultiplyNTT(ch, a)) for a in s1h]
    cs2  = [NTTInverse(MultiplyNTT(ch, a)) for a in s2h]
    z    = [Rq(y[i]) + cs1[i] for i in range(self.ell)]
    wcs2 = [w[i] - cs2[i] for i in range(self.k)]
    r0   = LowBitsVec(wcs2, self.gamma2)

    kappa += self.ell
    if inf_norm_vec(z) >= self.gamma1 - self.beta:      # unsafe: retry
        continue
    if inf_norm_ints(r0) >= self.gamma2 - self.beta:
        continue

    ct0 = [NTTInverse(MultiplyNTT(ch, a)) for a in t0h]
    h   = MakeHintVec([-p for p in ct0],
                      [wcs2[i] + ct0[i] for i in range(self.k)],
                      self.gamma2)
    if inf_norm_vec(ct0) >= self.gamma2:
        continue
    if sum(sum(p) for p in h) > self.omega:
        continue

    return self.sigEncode(ctilde, [centered(zi) for zi in z], h)
\end{lstlisting}

每个 \texttt{continue} 都对应图~\ref{fig:mldsa-abort} 中一次被丢弃的尝试；在大量签名上统计它们，就能重现参数表中的“重复次数”一行，对 ML-DSA-44 而言大约为四次。另请注意，$\kappa$ 在检查之前就已递增，因此每次重试都会抽取一个真正新鲜的掩码。

\paragraph{对照代数检验变换。}
\sagefile{fips204\_mldsa.sage} 末尾的函数 \texttt{verify\_ntt\_is\_a\_ring\_isomorphism()}
——与第~\ref{ch:fips203}~章中 ML-KEM 的检查同名，但作用于另一个环——借助 Sage
确认：$\mathrm{NTT}^{-1}$ 是 $\mathrm{NTT}$ 的逆；第 $j$ 个 NTT
系数是取值 $w(\zeta^{2\,\mathrm{BitRev8}(j)+1})$，这正是使 $T_q$ 成为直积的原因；以及 \textsc{MultiplyNTT} 如实传递了 Sage 所计算的
$R_q$ 中的乘法。它还检查
\textsc{MontgomeryReduce} 确实返回 $a\cdot 2^{-32} \bmod q$。

\paragraph{验证。}
该实现以 NIST 针对 FIPS~204 的 ACVP 向量~\cite{acvp}进行检验——包括密钥生成、签名生成与签名验证——并逐字节重现了这些向量。签名生成向量要求最高，因为它们覆盖了标准所允许的每一种组合：确定性签名与对冲（hedged）签名，纯消息编码与预哈希消息编码（遍及全部十二个获批的预哈希函数），外部与内部签名接口，以及外部 $\mu$ 模式——在该模式下，消息代表值在一个单独的密码模块中计算。在
\texttt{latex/} 目录下运行 \texttt{make kat} 即可。

\section{常见陷阱}

\paragraph{陷阱 1：以为掩码可以重复使用或取得很小。}
掩码 $\mathbf{y}$ 必须在每次尝试中都是新鲜的，并且足够大以隐藏 $c\,\mathbf{s}_1$。重复使用掩码，或为了避免拒绝而缩小掩码，都会泄露秘密。这正是设置中止循环的全部原因。

\paragraph{陷阱 2：以为中止循环会造成关于秘密的计时泄露。}
迭代次数取决于掩码，而掩码是新鲜的随机性，并不以可被利用的方式依赖于秘密。FIPS~204 由 $K\,\|\,rnd\,\|\,\mu$ 派生掩码种子：选择全零的 $rnd$ 得到确定性签名，而新鲜的 $rnd$ 则支持随机化签名或对冲签名。

\paragraph{陷阱 3：以为 $A$ 存储在密钥中。}
与 ML-KEM 一样，$A$ 每次都由 $\rho$ 重新生成。公钥只有 $(\rho,\mathbf{t}_1)$。

\section{参数集}

\begin{table}[H]
\centering
\begin{tabular}{lccc}
\toprule
方案 & $(k,\ell)$ & NIST 级别 & 相当于 \\
\midrule
ML-DSA-44 & $(4,4)$ & 2 & AES-128 \\
ML-DSA-65 & $(6,5)$ & 3 & AES-192 \\
ML-DSA-87 & $(8,7)$ & 5 & AES-256 \\
\bottomrule
\end{tabular}
\caption{三个标准化的 ML-DSA 参数集。它们都使用 $n=256$ 和 $q=8380417$。}
\end{table}

\section{小结}

ML-DSA 复用了 Kyber 的模格体系，并加入了一个真正全新的要素——带中止的 Fiat--Shamir：

\begin{itemize}
  \item 密钥生成是由单个种子 $\xi$ 进行的确定性扩展，与 FIPS 203 完全一样；
  \item 签名是一段借助哈希实现非交互化的承诺—挑战—响应记录；
  \item 响应 $\mathbf{z}=\mathbf{y}+c\,\mathbf{s}_1$ 携带秘密，因此只有当拒绝采样确认它完全隐藏了秘密时才会被发布；
  \item 不可伪造性依赖于模 LWE（隐藏密钥）和模 SIS（阻止短向量伪造）。
\end{itemize}

下一章仍停留在格上，但走了一条截然不同的签名路线：FN-DSA（Falcon）。它放弃 Fiat--Shamir，转而在 NTRU 格上采用哈希后签名（hash-and-sign）构造，并以此换来了三个标准中最紧凑的签名。

\printbibliography[heading=chapterrefs]
